\markboth{Abstract}{Abstract}
\section*{Abstract}

Since the crisis of 2021-2022, the European natural gas market has become highly sensitive to supply disruptions, which can cause extreme price volatility. Widely used option pricing models such as Black-76 assume log returns that are normally distributed with constant volatility. However, observed gas prices show heavy tails, asymmetry, and volatility clustering and these stylized facts are not consistent with the assumptions of Black-76. Generative adversarial networks (GANs) learn the distribution of the data directly from observations. This thesis applies the conditional GAN frameworks ForGAN and Fin-GAN to the daily Dutch Title Transfer Facility (TTF) futures returns. A grid of ten cost function variants and eleven network and training configurations is evaluated on three stylized facts. The most reproducible model is selected and generates price paths for twenty valuation dates in the test period. European options are priced by Monte Carlo simulation. The resulting implied volatilities are compared with the market quotes and with three benchmark models. The selected model reproduces the stylized facts in part, but the option prices it generates stay below the market and below all three benchmarks. The spread of the simulated returns correlates \num{0.752} with the realized volatility in the condition window. No benchmark shows this dependence on the market state. The conditioning works, while the level of the generated volatility remains for future work.